This section reviews deep learning–based motor imagery EEG classification, covering general deep learning concepts, MI-specific challenges, end-to-end CNN architectures, and the research position of this study.

\subsection{Deep Learning Concept}

Feedforward neural networks (FNNs) underpin modern deep learning but lack structural awareness, typically serving as final decision layers for structured data like EEG~\cite{Walther2023A}. For feature extraction, Convolutional networks (CNNs) capture spatial patterns via local connectivity~\cite{Shiri2023A}, while Recurrent networks (LSTMs) manage temporal dependencies through internal memory states~\cite{Mienye2024Recurrent}. More recently, Transformers have advanced sequence modeling by utilizing self-attention mechanisms to capture global relationships in parallel~\cite{Hasan2023MALS-Net}. To ensure robust performance, training relies on optimizers (e.g., Adam) for convergence~\cite{Bartlett2021Deep} and regularization techniques (e.g., dropout) to prevent overfitting and define generalization capacity~\cite{Poggio2019Theoretical,Moradi2019A}.

\subsection{Deep Learning for EEG Signal Classification}

The shift from traditional ML to deep learning in BCI mirrors broader trends in computer vision and speech processing. CNNs enable "end-to-end" learning, taking raw or minimally preprocessed EEG as input and directly outputting class probabilities~\cite{schirrmeister2017deep,elouahidi2024strong,Dose2018An}. This eliminates manual feature engineering steps—artifact removal, bandpass filtering, feature selection—common in traditional pipelines like FBCSP~\cite{Peng2024,lawhern2018eegnet,Tibrewal2022Classification}.

However, applying DL to EEG requires more than importing vision architectures. EEG has fundamentally different structure: anisotropic correlations across temporal and spatial dimensions, unlike isotropic image correlations~\cite{kollod2023deep,padfield2019eeg}. The temporal dimension encodes frequency information essential for SMR recognition, while the spatial dimension (channels) provides topographic source localization~\cite{Belwafi2020An,Wei2023Decoding}. Furthermore, BCI faces severe data scarcity—typically hundreds of trials per subject (e.g., 288 in BCI Competition IV 2a) versus millions of images in computer vision~\cite{BCI2008GrazDataset}, making standard deep networks prone to overfitting.

CNNs dominate EEG decoding due to architectural compatibility with EEG properties. Local connectivity and weight sharing enable efficient detection of patterns like waveform complexes or oscillatory bursts across entire time series~\cite{Zhang2021EEG-inception:}. Temporal convolutions function as learnable FIR filters, automatically discovering task- and subject-specific bandpass features in mu, beta, or gamma bands~\cite{2017LearningDR,Wu2025Learning-based,Fukumori2020Epileptic}. Spatial convolutions linearly combine electrode data as learnable spatial filters, analogous to CSP or beamforming, enhancing task-relevant sources while suppressing noise~\cite{Borra2020Interpretable}. This architecture offers interpretability: visualizing convolutional weights can confirm neurophysiologically plausible learning, such as spatial filters focusing on C3/C4 electrodes for hand motor imagery~\cite{lawhern2018eegnet}.

\subsection{Research Positions}
This study evaluates EEGNet and ShallowConvNet for motor imagery classification using data acquired via the low-cost OpenBCI Cyton board. Unlike standard benchmarks, this approach assesses model robustness under the realistic noise and variability of accessible hardware. Through identical processing protocols, the analysis determines the viability of compact CNN architectures for portable, real-world EEG applications.
